\section*{Capítulo \ref{ch:g12:stereochemistry} --- Estereoquímica: quiralidad e isómeros}

\begin{solution}{exo:g12:stereochemistry:1}
Butan-2-ol: el carbono 2 (\ce{H}, \ce{OH}, \ce{CH3}, \ce{CH2CH3}).
Propan-2-ol: ninguno (el carbono 2 lleva dos \ce{CH3} idénticos).
2-Clorobutano: el carbono 2. Ácido láctico: el carbono 2 (\ce{H}, \ce{OH},
\ce{CH3}, \ce{COOH}).
\end{solution}

\begin{solution}{exo:g12:stereochemistry:2}
Primero: los dos \ce{Cl} (prioritarios frente a \ce{H}) del mismo lado: \omterm{def:g12:stereochemistry:z-e}{Z}.
Segundo: a la izquierda, el \ce{CH3} es prioritario, en el lado de arriba;
a la derecha, el \ce{CH2CH3} es prioritario, en el lado de abajo: \omterm{def:g12:stereochemistry:z-e}{E}.
\end{solution}

\begin{solution}{exo:g12:stereochemistry:3}
\omterm{def:g12:stereochemistry:chiral}{Quirales}: el guante y el tornillo. No \omterm{def:g12:stereochemistry:chiral}{quirales}: la cuchara y el cubo
(cada uno tiene un plano de simetría). El \ce{CHFClBr} es \omterm{def:g12:stereochemistry:chiral}{quiral} (cuatro
\omterm{def:g7:atoms-and-molecules:atom}{átomos} distintos sobre el carbono); el \ce{CH2Cl2} no lo es.
\end{solution}

\begin{solution}{exo:g12:stereochemistry:4}
Una \omterm{def:g6:pure-substances-and-mixtures:mixture}{mezcla} de cantidades iguales de dos \omterm{def:g12:stereochemistry:enantiomer}{enantiómeros}. La carvona
racémica contiene los dos \omterm{def:g12:stereochemistry:enantiomer}{enantiómeros}, así que la nariz percibe los dos
olores a la vez.
\end{solution}

\begin{solution}{exo:g12:stereochemistry:5}
No: el carbono 1 lleva dos \omterm{def:g7:atoms-and-molecules:atom}{átomos} idénticos, \ce{H} y \ce{H};
intercambiarlos da la misma \omterm{def:g7:atoms-and-molecules:molecule}{molécula}.
\end{solution}

\begin{solution}{exo:g12:stereochemistry:6}
\setchemfig{atom sep=2em}
\chemfig{C(-[2]CH_2CH_3)(-[4]H_3C)(<[:-30]OH)(<:[:-150]H)} \quad espejo
\quad \chemfig{C(-[2]CH_2CH_3)(-[0]CH_3)(<[:-150]HO)(<:[:-30]H)}
\end{solution}

\begin{solution}{exo:g12:stereochemistry:7}
2-Clorobutano: 2 (un \omterm{def:g12:stereochemistry:chiral}{carbono asimétrico}). 2,3-Diclorobutano: 3 (dos
\omterm{def:g12:stereochemistry:chiral}{carbonos asimétricos} con mitades iguales: una pareja de \omterm{def:g12:stereochemistry:enantiomer}{enantiómeros} y
una forma meso). Pent-2-eno: 2 (\omterm{def:g12:stereochemistry:z-e}{Z} y \omterm{def:g12:stereochemistry:z-e}{E}).
\end{solution}

\begin{solution}{exo:g12:stereochemistry:8}
\omterm{def:g12:stereochemistry:enantiomer}{Enantiómeros}; \omterm{def:g12:stereochemistry:diastereomer}{diastereoisómeros}; \omterm{def:g12:stereochemistry:diastereomer}{diastereoisómeros}; la misma \omterm{def:g7:atoms-and-molecules:molecule}{molécula}
(el ácido meso-tartárico es superponible a su imagen especular).
\end{solution}

\begin{solution}{exo:g12:stereochemistry:9}
En la \omterm{def:g12:stereochemistry:conformation}{conformación} alternada con los grupos \ce{CH3} opuestos, los
grupos grandes están lo más alejados posible: es la más estable. En la
eclipsada, los dos \ce{CH3} están enfrentados y se estorban.
\end{solution}

\begin{solution}{exo:g12:stereochemistry:10}
Sus dos mitades, cada una \ce{CH(OH)COOH}, son imagen especular una de
la otra: un plano entre los dos carbonos centrales corta la \omterm{def:g7:atoms-and-molecules:molecule}{molécula} en
dos mitades especulares. Una \omterm{def:g7:atoms-and-molecules:molecule}{molécula} con un plano de simetría es su
propia imagen especular: no es \omterm{def:g12:stereochemistry:chiral}{quiral}.
\end{solution}

\begin{solution}{exo:g12:stereochemistry:11}
\setchemfig{atom sep=1.7em}
(\omterm{def:g12:stereochemistry:z-e}{Z})-pent-2-eno \chemfig{C(-[:150]H_3C)(-[:210]H)=C(-[:30]CH_2-CH_3)(-[:-30]H)}
y (\omterm{def:g12:stereochemistry:z-e}{E})-pent-2-eno \chemfig{C(-[:150]H_3C)(-[:210]H)=C(-[:30]H)(-[:-30]CH_2-CH_3)}.
\end{solution}

\begin{solution}{exo:g12:stereochemistry:12}
Como para el ácido tartárico, con \ce{Br} en lugar de \ce{OH} y \ce{CH3}
en lugar de \ce{COOH}: los dos \ce{Br} con cuña llena, los dos con cuña
de trazos (una pareja de \omterm{def:g12:stereochemistry:enantiomer}{enantiómeros}), y uno con cuña llena y otro con
cuña de trazos, que tiene un plano de simetría: la forma meso. Tres
\omterm{def:g12:stereochemistry:stereoisomer}{estereoisómeros}.
\end{solution}

\begin{solution}{exo:g12:stereochemistry:13}
La mitad. El químico puede separar los dos \omterm{def:g12:stereochemistry:enantiomer}{enantiómeros} (como hizo
Pasteur, por otros medios), o fabricar solo el \omterm{def:g12:stereochemistry:enantiomer}{enantiómero} útil, con una
reacción que sea ella misma \omterm{def:g12:stereochemistry:chiral}{quiral} (un catalizador \omterm{def:g12:stereochemistry:chiral}{quiral}, una enzima,
un material de partida \omterm{def:g12:stereochemistry:chiral}{quiral}).
\end{solution}

\begin{solution}{exo:g12:stereochemistry:14}
Tres: (2E,4E), (2Z,4Z) y (2E,4Z); este último es la misma \omterm{def:g7:atoms-and-molecules:molecule}{molécula} que
el (2Z,4E) leído desde el otro extremo.
\end{solution}

\begin{solution}{exo:g12:stereochemistry:15}
$2^3 = 8$ \omterm{def:g12:stereochemistry:stereoisomer}{estereoisómeros}, es decir, 4 parejas de \omterm{def:g12:stereochemistry:enantiomer}{enantiómeros}. Cada
\omterm{def:g12:stereochemistry:stereoisomer}{estereoisómero} tiene un \omterm{def:g12:stereochemistry:enantiomer}{enantiómero} y $8 - 2 = 6$ \omterm{def:g12:stereochemistry:diastereomer}{diastereoisómeros}.
\end{solution}

\begin{solution}{pb:g12:stereochemistry:1}
\textbf{1.} \ce{C4H6O6}; $M = 4 \times 12.0 + 6 \times 1.0 + 6 \times 16.0
= \qty{150.0}{g/mol}$.

\textbf{2.} Dos grupos carboxilo (ácido) y dos grupos hidroxilo
(\omterm{def:g11:functional-groups:alcohol}{alcohol}).

\textbf{3.} Los dos carbonos centrales, cada uno unido a \ce{H}, \ce{OH},
\ce{COOH} y \ce{CH(OH)COOH}.

\textbf{4.} $2^2 = 4$.

\textbf{5.} Los dos con cuñas llenas.

\textbf{6.} Los dos \ce{OH} con cuñas de trazos; no se puede superponer
al L: es el ácido D-tartárico, su \omterm{def:g12:stereochemistry:enantiomer}{enantiómero}.

\textbf{7.} Girada en torno al enlace central de modo que los dos
\ce{OH} miren hacia el mismo lado, la \omterm{def:g7:atoms-and-molecules:molecule}{molécula} tiene un plano entre sus
dos carbonos centrales que refleja una mitad sobre la otra.

\textbf{8.} No: es el ácido meso-tartárico.

\textbf{9.} De las cuatro combinaciones (llena, llena), (trazos, trazos),
(llena, trazos), (trazos, llena), las dos últimas son la misma \omterm{def:g7:atoms-and-molecules:molecule}{molécula},
la forma meso, porque las dos mitades de la \omterm{def:g7:atoms-and-molecules:molecule}{molécula} son iguales.

\textbf{10.} L y D: \omterm{def:g12:stereochemistry:enantiomer}{enantiómeros}. L y meso: \omterm{def:g12:stereochemistry:diastereomer}{diastereoisómeros}.

\textbf{11.} También a \qty{169}{\celsius}: los \omterm{def:g12:stereochemistry:enantiomer}{enantiómeros} tienen las
mismas propiedades físicas.

\textbf{12.} No: es un \omterm{def:g12:stereochemistry:diastereomer}{diastereoisómero} del L, un compuesto distinto con
sus propias propiedades.

\textbf{13.} Los receptores de la nariz son \omterm{def:g12:stereochemistry:chiral}{quirales} y distinguen las
imágenes especulares; un termómetro, que no es \omterm{def:g12:stereochemistry:chiral}{quiral}, no puede.

\textbf{14.} No, es una \omterm{def:g6:pure-substances-and-mixtures:mixture}{mezcla} de dos compuestos; en conjunto no es
ópticamente activa, porque los dos \omterm{def:g12:stereochemistry:enantiomer}{enantiómeros} se compensan.

\textbf{15.} Un solo \omterm{def:g12:stereochemistry:enantiomer}{enantiómero}: cada cristal estaba formado solo por L
o solo por D.

\textbf{16.} \qty{5.0}{g} de cada uno.

\textbf{17.} La forma meso es un compuesto distinto, que no forma parte
de la \omterm{def:g12:stereochemistry:enantiomer}{mezcla racémica} de L y D.

\textbf{18.} \omterm{def:g2:mixing-and-dissolving:dissolve}{Disolviendo} cada montón: las disoluciones desvían la luz
polarizada en la misma medida y en sentidos opuestos (y los cristales
son imágenes especulares).

\textbf{19.} Cuatro: sus dos \omterm{def:g12:stereochemistry:chiral}{carbonos asimétricos} llevan grupos
distintos (\ce{Cl} en uno, \ce{OH} en el otro), así que ninguna
combinación tiene un plano de simetría y no hay dos combinaciones que
sean la misma \omterm{def:g7:atoms-and-molecules:molecule}{molécula}.

\textbf{20.} Tres: L, D y meso.
\end{solution}
